% Generated by scientific-research-assistant. Compile with:
%   pdflatex report && bibtex report && pdflatex report && pdflatex report
% Use xelatex or lualatex instead if titles contain characters pdflatex rejects.
\documentclass[11pt,a4paper]{article}
\usepackage[utf8]{inputenc}
\usepackage[T1]{fontenc}
\usepackage[margin=2.5cm]{geometry}
\usepackage{booktabs,longtable,array}
\usepackage[hidelinks]{hyperref}
\usepackage{url}
\setlength{\parindent}{0pt}
\setlength{\parskip}{0.6em}

\title{Adversarial audit: do syntax-tree GCNs help ABSA when the encoder is held fixed?}
\date{2026-10-18}

\begin{document}
\maketitle


\textbf{Report type:} Research Opportunity Report\\
\textbf{Generated:} 2026-10-18

\textbf{Evidence tags:} [V] verified against a cited source or API response; [E] reasoned estimate; [S] speculative suggestion.

{\small
\begin{longtable}{>{\raggedright\arraybackslash}p{0.307\textwidth}>{\raggedright\arraybackslash}p{0.307\textwidth}>{\raggedright\arraybackslash}p{0.307\textwidth}}
\toprule
\textbf{Constraint} & \textbf{Value} & \textbf{Origin} \\
\midrule
\endhead
field & Aspect-based sentiment analysis with dependency-graph neural networks & stated by user \\
experience & undergraduate, comfortable with Python & default (not stated by user) \\
hardware & one consumer GPU, 8 GB VRAM & default (not stated by user) \\
gpu vram gb & 8 & default (not stated by user) \\
time budget weeks & 8 & default (not stated by user) \\
language & Python & default (not stated by user) \\
difficulty & moderate & default (not stated by user) \\
mode & any & default (not stated by user) \\
\bottomrule
\end{longtable}
}

\section{Summary}

The question as posed is already answered: Dai et al. ran RoBERTa with and without ASGCN, PWCN and RGAT tree heads on the same SemEval data and found no significant gain \cite{dai2021syntax}. What remains open is whether that null result, and the 2023-2026 claims of gains, survive a sound protocol: the closest codebases select models on the test set, a selection habit worth 1-2 points elsewhere in ABSA \cite{mukherjee2021repro}, about the size of every claimed syntax gain. A protocol-corrected replication with structure controls is a conditional go; a novel framing (frozen vs fine-tuned encoder) is plausible but unverified.

\section{Decisions}

Each shortlisted research question gets go, conditional go or no-go. A go requires full-text reading of the closest work and an explicit check for leakage, weak baselines and evaluation flaws; anything short of that is at most a conditional go.

{\small
\begin{longtable}{>{\raggedright\arraybackslash}p{0.184\textwidth}>{\raggedright\arraybackslash}p{0.184\textwidth}>{\raggedright\arraybackslash}p{0.184\textwidth}>{\raggedright\arraybackslash}p{0.184\textwidth}>{\raggedright\arraybackslash}p{0.184\textwidth}}
\toprule
\textbf{Research question} & \textbf{Decision} & \textbf{Novelty confidence} & \textbf{Mode} & \textbf{Main conditions} \\
\midrule
\endhead
As stated: do syntax-tree GCNs help when the pretrained encoder is held fixed? & NO-GO & very\_low & original & not reported \\
D1 (replication): protocol-corrected re-test of Dai et al.'s RoBERTa tree-head comparison & CONDITIONAL GO & low & replication & Read Lebmeier et al. through the library; if it already re-runs a syntax-GCN model with validation splits and uncertainty, narrow to the DeBERTa-v3 part or drop D1.; Fix the validation protocol, seeds and the smallest effect of interest (e. \\
D2 (extension): is it syntax, or just an extra graph head? Structure controls on a fixed fine-tuned encoder & CONDITIONAL GO & low & extension & Run D1 first so the baseline and protocol are trusted.; Screen ablation tables of the 5-10 most cited 2022-2026 syntax-GCN ABSA papers for random or chain-tree controls on PLMs. \\
D3 (novel investigation): frozen versus fine-tuned encoder - does syntax matter only when the encoder cannot adapt? & CONDITIONAL GO & low & original & Run a dedicated challenge and web search on frozen-encoder syntax probing (e.g. 'frozen BERT features dependency GCN') before committing. \\
D4 (extension): parser quality as the moderator of syntax gains in ABSA & NO-GO & low & extension & not reported \\
\bottomrule
\end{longtable}
}

\section{What was verified, and how}

{\small
\begin{longtable}{>{\raggedright\arraybackslash}p{0.307\textwidth}>{\raggedright\arraybackslash}p{0.307\textwidth}>{\raggedright\arraybackslash}p{0.307\textwidth}}
\toprule
\textbf{Dimension} & \textbf{Status} & \textbf{Meaning} \\
\midrule
\endhead
Source existence (verify-refs) & verified: 13 & The cited record exists and its title matches. Says nothing about what the paper shows. \\
Claim support & papers read: abstract: 5, full text: 6, title only: 2; claims: estimate: 7, speculative: 2, verified: 35 & Verified claims were checked against their source at the read level shown. \\
Novelty confidence & As stated: do syntax-tree GCNs help when the pretrained encoder is held fixed?: very\_low (search ceiling low); D1 (replication): protocol-corrected re-test of Dai et al.'s RoBERTa tree-head comparison: low (search ceiling low); D2 (extension): is it syntax, or just an extra graph head? Structure controls on a fixed fine-tuned encoder: low (search ceiling low); D3 (novel investigation): frozen versus fine-tuned encoder - does syntax matter only when the encoder cannot adapt?: low (search ceiling low); D4 (extension): parser quality as the moderator of syntax gains in ABSA: low (search ceiling low) & How far the searches support a gap; never above medium. \\
Reproducibility & nothing executed & Whether results were re-run, not whether code exists. \\
\bottomrule
\end{longtable}
}

\section{How the literature was searched}

\begin{itemize}
\item \textbf{Sources that answered:} openalex, arxiv, crossref, huggingface
\item \textbf{Sources that failed or were skipped:} semantic\_scholar (rate\_limited)
\item \textbf{Search strategies:} challenge with design elements and seed (arXiv 2104.04986), citation expansion, web search, full-text reading (fulltext command), code reading, data leakage check
\item \textbf{Date range:} unrestricted
\item \textbf{Records retrieved / after deduplication:} 372 / None
\item \textbf{Search date:} 2026-10-18
\item \textbf{Queries:} \texttt{does syntax help pretrained transformers aspect-based sentiment}; \texttt{dependency tree graph neural network on top of BERT ablation sentiment}; \texttt{syntactic information unnecessary strong RoBERTa baseline aspect sentiment}; \texttt{syntax-agnostic versus syntax-aware models frozen encoder probing}; \texttt{random trees versus dependency trees graph convolution evaluation}; \texttt{web: does syntax help pretrained transformers dependency tree GNN ablation aspect-based sentiment random trees}; \texttt{web: aspect-based sentiment GCN reproducibility seeds variance same encoder syntax benefit}; \texttt{web: aspect sentiment random tree / fully connected graph replace dependency tree GCN ablation BERT}; \texttt{web: empirical study syntax ABSA pre-trained language models revisiting}; \texttt{web: ABSA evaluation model selection on test set no development set reproducibility}
\item Reused from the NLP+GNN opportunity report: candidate O2, its challenge (work/challenge\_absa.json), references and dataset checks.
\item The challenge returned no\_close\_match\_in\_this\_search and flagged an unrelated pharmacy paper on the design element 'frozen pretrained encoder'; Sachan et al., Mukherjee et al. and Lebmeier et al. were found by web search.
\item Full texts read via arXiv source: 2104.04986, 2008.09084, 2101.09449, 2110.00171, 2103.04350 (targeted), 2311.04467 (targeted). Lebmeier et al. is behind a publisher login (institutional copy is a viewer link).
\item Code read: RoBERTaABSA (finetune.py, ASGCN/train.py, data\_utils.py, args and logs), GeneZC/ASGCN train.py, shenwzh3/RGAT-ABSA run.py and trainer.py.
\end{itemize}

\section{What has already been done}

{\small
\begin{longtable}{>{\raggedright\arraybackslash}p{0.115\textwidth}>{\raggedright\arraybackslash}p{0.115\textwidth}>{\raggedright\arraybackslash}p{0.115\textwidth}>{\raggedright\arraybackslash}p{0.115\textwidth}>{\raggedright\arraybackslash}p{0.115\textwidth}>{\raggedright\arraybackslash}p{0.115\textwidth}>{\raggedright\arraybackslash}p{0.115\textwidth}>{\raggedright\arraybackslash}p{0.115\textwidth}}
\toprule
\textbf{Paper} & \textbf{Published} & \textbf{Method} & \textbf{Datasets} & \textbf{Findings} & \textbf{Stated limitations} & \textbf{Read} & \textbf{Link} \\
\midrule
\endhead
Does syntax matter? A strong baseline for Aspect-based Sentiment Analysis with RoBERTa (Dai et al. 2021) & 2021 & Compares parser trees, PTM-induced trees, FT-PTM-induced trees and left/right-chain trees inside ASGCN, PWCN and RGAT; RoBERTa with an MLP head as baseline & Rest14, Laptop14, Twitter; three SemEval-2016 non-English sets & Full text: tree-structure comparison uses GloVe embeddings, 3 runs. SOTA table: RoBERTa-MLP 87.37/83.78 accuracy (Rest14/Laptop14) vs RoBERTa+ASGCN, PWCN, RGAT with parser or FT-RoBERTa trees between 86.87 and 87.52 (Rest14) and 81.66-84.16 (Laptop14); text: 'incorporating trees over the RoBERTa brings no significant improvement'. No variance reported for that table. & Code (README): released logs are single runs; finetune.py uses the test set as dev for model selection; released tree-model code uses GloVe only, so the RoBERTa+tree rows cannot be re-run from the repository; no licence & full text & \href{https://doi.org/10.18653/v1/2021.naacl-main.146}{link} \\
Do Syntax Trees Help Pre-trained Transformers Extract Information? (Sachan et al. 2020) & 2020-08-20 & Syntax-GNN over BERT outputs (late fusion) and syntax inside attention (joint fusion); same encoder with and without dependency trees; gold vs in-domain vs off-the-shelf parses & CoNLL-2005/2012 SRL, OntoNotes NER, TACRED-style RE & Full text: 5 runs; gains on SRL (>2.5 F1) only with gold parses, contingent on parse quality; no gain on NER (within one standard deviation of bert-base); late-fusion gains of about 2 F1 also with RoBERTa when gold parses are used & Information extraction, not sentiment; encoder always fine-tuned & full text & \href{https://arxiv.org/abs/2008.09084}{link} \\
Reproducibility, Replicability and Beyond: Assessing Production Readiness of Aspect Based Sentiment Analysis in the Wild (Mukherjee et al. 2021) & 2021 & Re-runs six ABSA models (memory-network and BERT-based) with test-set model selection vs a 15\% held-out split from training & SemEval-2014 Laptop and Restaurant; two new e-commerce datasets & Full text: reproduced accuracy 1-2 points below reported; a further 1-2 point drop when model selection uses a held-out split instead of the test set; 5 runs & No syntax or GCN models among the six & full text & \href{https://doi.org/10.1007/978-3-030-72240-1_7}{link} \\
On the Current State of Reproducibility and Reporting of Uncertainty for Aspect-Based Sentiment Analysis (Lebmeier et al. 2023) & 2023-01-01 & Holistic comparison of six (near) state-of-the-art transformer-based ABSA architectures with fixed preprocessing, train/validation splits, metrics and uncertainty quantification & Five public benchmark datasets & Abstract: reported results are hardly reproducible; results vary notably across data splits, so uncertainty must be reported & Full text behind publisher login; the six architectures are not named in the abstract & abstract & \href{https://doi.org/10.1007/978-3-031-26390-3_31}{link} \\
BERT4GCN: Using BERT Intermediate Layers to Augment GCN for Aspect-based Sentiment Classification (Xiao et al. 2021) & 2021 & GCN over dependency graphs supplemented with BERT intermediate-layer features and relative positions & Twitter, Laptop14, Rest14 & Full text: ablation shows GCN over BERT sequential features without the supplemented dependency graph already reaches state-of-the-art and the graph adds 'just slightly'; its RoBERTa-MLP baseline is 81.11 on Laptop14 (Dai et al. report 83.78); averages over 10 runs ('10-fold') & Model-selection protocol not stated in the text read & full text & \href{https://doi.org/10.18653/v1/2021.emnlp-main.724}{link} \\
Syntax-BERT: Improving Pre-trained Transformers with Syntax Trees (Bai et al. 2021) & 2021 & Syntax-tree masks inside pre-trained transformers; ablation replaces syntactic masks with random masks of equal parameter size & GLUE-style sentence classification tasks & Full text (targeted passages): a random decomposition of the network also gives moderate improvements & Not ABSA; random-mask numbers not extracted here & full text & \href{https://doi.org/10.18653/v1/2021.eacl-main.262}{link} \\
RDGCN: Reinforced Dependency Graph Convolutional Network for Aspect-based Sentiment Analysis (Zhao et al. 2023) & 2023-11-08 & Reinforced dependency GCN with RL-tuned distance-importance function & not reported & Full text (targeted passages): 'Since the datasets do not contain the validation set, the test accuracy is used to implement the reward function' & Test labels enter training through the RL reward & full text & \href{https://arxiv.org/abs/2311.04467}{link} \\
Selective Attention Based Graph Convolutional Networks for Aspect-Level Sentiment Classification (Hou et al. 2021) & 2021 & Selective-attention GCN over dependency trees on BERT & not reported & Web-search snippet (not verified against text): BERT baselines with 2- and 4-layer GCN, 7-run averages & arXiv source contains only the supplement; main text not read & abstract & \href{https://doi.org/10.18653/v1/2021.textgraphs-1.8}{link} \\
Aspect-based Sentiment Classification with Aspect-specific Graph Convolutional Networks (Zhang et al. 2019) & 2019 & Aspect-specific GCN over dependency trees (GloVe/LSTM) & not reported & not reported & not reported & abstract & \href{https://doi.org/10.18653/v1/d19-1464}{link} \\
Relational Graph Attention Network for Aspect-based Sentiment Analysis (Wang et al. 2020) & 2020 & Relational graph attention over aspect-oriented dependency trees & not reported & not reported & Official code (MIT) selects the best epoch on the test set (run.py, trainer.py) & abstract & \href{https://doi.org/10.18653/v1/2020.acl-main.295}{link} \\
FITE-GAT: Enhancing aspect-level sentiment classification with FT-RoBERTa induced trees and graph attention network (Fan et al. 2024) & 2024-11-23 & Title: FT-RoBERTa induced trees with graph attention; abstract not retrieved & not reported & not reported & not reported & title only & \href{https://doi.org/10.1016/j.eswa.2024.125890}{link} \\
Syntax-Enhanced Pretrained Language Models for Aspect-Level Sentiment Classification (Li et al. 2024) & 2024-12-24 & SE-BERT: syntax-enhanced transformer layers replacing BERT's final layers & not reported & Abstract states that existing methods improve aspect-level sentiment classification by adding syntactic GCN layers on BERT & not reported & abstract & \href{https://doi.org/10.1109/tcss.2024.3514901}{link} \\
Multi-View Attention Syntactic Enhanced Graph Convolutional Network for Aspect-based Sentiment Analysis (Huang et al. 2025) & 2025-01-27 & Multi-view attention syntactic enhanced GCN for ABSA & not reported & Example of continued dependency-GCN ABSA work after 2021 & not reported & title only & \href{https://arxiv.org/abs/2501.15968}{link} \\
\bottomrule
\end{longtable}
}

Method, dataset, finding and limitation cells summarise the cited paper's abstract or text; 'not reported' means the information was not found there.

\subsection{Directions that are already well explored}

\begin{itemize}
\item [E] New dependency-GCN architectures on SemEval-2014 (many 2022-2026 papers; titles screened) \cite{li2024sebert,zhao2023rdgcn}
\item [V] On Rest14 and Laptop14, RoBERTa+tree heads lie within about +/-1.1 (Rest14) and -2.1 to +0.4 (Laptop14) accuracy points of RoBERTa-MLP, with no variance reported. \cite{dai2021syntax}
\item [V] Whether syntax helps fine-tuned transformers on information extraction \cite{sachan2021syntax}
\item [V] Reproducibility of ABSA results in general \cite{mukherjee2021repro}
\end{itemize}

\section{Ranked research opportunities}

Each opportunity below survived an explicit search for work that already addresses it. None is claimed to be novel: the novelty confidence states how far the search supports the gap, and what it could have missed.

\subsection{1. D1 (replication): protocol-corrected re-test of Dai et al.'s RoBERTa tree-head comparison}
\textbf{Type of gap:} Limitation reported in prior work\\
\textbf{Recommended mode:} replication\\
\textbf{Revised:} Narrowed from the original question after the full texts showed the comparison exists but was run with test-set model selection, without variance and without released code for the RoBERTa+tree rows.

\textbf{Research question.} On the shipped SemEval-2014 Laptop and Restaurant splits, does the reported null result (no gain from tree heads on RoBERTa) survive validation-based model selection, a tuned baseline and 10+ seeds, and does it hold for DeBERTa-v3-base?

\textbf{Motivation.} The headline null result is cited as settled but rests on test-selected, single-table numbers \cite{dai2021syntax}.

\textbf{Supporting papers.} \cite{dai2021syntax,mukherjee2021repro,lebmeier2023repro}

\textbf{Closest existing work.}
\begin{itemize}
\item \cite{dai2021syntax} [read: full text; overlap: same experiment]: The study being replicated; the replication changes the protocol, not the question.
\item \cite{lebmeier2023repro} [read: abstract; overlap: unknown]: Re-runs six transformer ABSA architectures with validation splits and uncertainty; whether any is a syntax-GCN model is unknown.
\item \cite{mukherjee2021repro} [read: full text; overlap: partial]: Quantifies test-selection inflation, but for memory-network and BERT models, not tree heads.
\end{itemize}

\textbf{Evidence for the gap.}
\begin{itemize}
\item [V] On Rest14 and Laptop14, RoBERTa+tree heads lie within about +/-1.1 (Rest14) and -2.1 to +0.4 (Laptop14) accuracy points of RoBERTa-MLP, with no variance reported. \cite{dai2021syntax}
\item [V] Official code of the closest studies selects on the test set: RoBERTaABSA finetune.py passes the test split as dev\_data and reloads the best model; ASGCN (original and Dai et al.'s copy) keeps max\_test\_acc and max\_test\_f1 separately, so accuracy and F1 can come from different epochs; RGAT run.py reports the best evaluation on the test set. \cite{dai2021syntax,zhang2019asgcn,wang2020rgat}
\item [V] Selecting on a 15\% held-out split instead of the test set costs a further 1-2 accuracy points for six non-syntax ABSA models, about the size of the claimed syntax gains. \cite{mukherjee2021repro}
\item [V] The same RoBERTa-MLP baseline is reported at 83.78 (Dai et al.) and 81.11 (BERT4GCN) on Laptop14, so cross-paper syntax gains can be baseline artefacts. \cite{dai2021syntax,xiao2021bert4gcn}
\end{itemize}

\textbf{Potential contribution.} A checked, variance-aware answer to a question the field treats as settled; either outcome is reportable.

\textbf{Novelty confidence:} low (scale: very\_low, low, medium; a search cannot establish more). Replication value, not novelty. Lebmeier et al.'s architectures are unknown; a paywalled or non-English replication could exist.
Attempts to disprove novelty: 10 search(es); signal \texttt{no\_close\_match\_in\_this\_search}. Queries: \texttt{does syntax help pretrained transformers aspect-based sentiment}; \texttt{dependency tree graph neural network on top of BERT ablation sentiment}; \texttt{syntactic information unnecessary strong RoBERTa baseline aspect sentiment}; \texttt{syntax-agnostic versus syntax-aware models frozen encoder probing}; \texttt{random trees versus dependency trees graph convolution evaluation}; \texttt{web: does syntax help pretrained transformers dependency tree GNN ablation aspect-based sentiment random trees}; \texttt{web: aspect-based sentiment GCN reproducibility seeds variance same encoder syntax benefit}; \texttt{web: aspect sentiment random tree / fully connected graph replace dependency tree GCN ablation BERT}

\textbf{Feasibility.}
\begin{itemize}
\item [V] Shipped splits and pre-parsed dependency trees (head, deprel fields) are in the RoBERTaABSA repository; RoBERTa-base and DeBERTa-v3-base are MIT-licensed on Hugging Face.
\item [E] About 3,600 training aspects; base encoders at batch 16 fit 8 GB with mixed precision; roughly 2 datasets x 2 encoders x 3 heads x 10 seeds = 120 short runs, about 15-30 GPU-hours.
\item [V] Tree heads must be reimplemented on current transformers (repository pins transformers 2.11, PyTorch 1.5, fastNLP).
\end{itemize}

\textbf{Methodological risks.}

{\small
\begin{longtable}{>{\raggedright\arraybackslash}p{0.23\textwidth}>{\raggedright\arraybackslash}p{0.23\textwidth}>{\raggedright\arraybackslash}p{0.23\textwidth}>{\raggedright\arraybackslash}p{0.23\textwidth}}
\toprule
\textbf{Risk} & \textbf{Status} & \textbf{Evidence} & \textbf{Mitigation} \\
\midrule
\endhead
model selection on test & found & [V] Official code of the closest studies selects on the test set: RoBERTaABSA finetune.py passes the test split as dev\_data and reloads the best model; ASGCN (original and Dai et al.'s copy) keeps max\_test\_acc and max\_test\_f1 separately, so accuracy and F1 can come from different epochs; RGAT run.py reports the best evaluation on the test set. \cite{dai2021syntax,zhang2019asgcn,wang2020rgat} & Validation split carved from training, grouped by sentence; select checkpoints and hyperparameters on it only; touch the test set once. \\
metric definition & found & [V] ASGCN reports maximum accuracy and maximum F1 from possibly different epochs. \cite{zhang2019asgcn} & Report accuracy and macro-F1 of one selected checkpoint. \\
insufficient runs & found & [V] Dai et al.: 3 runs for the tree comparison, no variance for the RoBERTa table, released logs single-run; Lebmeier et al. report notable variance across data splits. \cite{dai2021syntax,lebmeier2023repro} & At least 10 seeds x repeated train/validation splits; paired comparisons; confidence intervals. \\
weak baseline & found & [V] The same RoBERTa-MLP baseline is reported at 83.78 (Dai et al.) and 81.11 (BERT4GCN) on Laptop14, so cross-paper syntax gains can be baseline artefacts. \cite{dai2021syntax,xiao2021bert4gcn} & Tune the plain-encoder baseline with the same budget and search space as the tree heads. \\
protocol mismatch & found & [V] No public code for the RoBERTa+tree rows; the Hugging Face SetFit copy has 654 laptop test rows vs 638 aspects in the shipped split. \cite{dai2021syntax} & Use the shipped split files; reimplement heads on one codebase; label the work a replication. \\
confound & not checked & Tree heads add parameters and a separate learning rate; Dai et al. note RoBERTa-based tree models are hard to optimise and need different learning rates. & Parameter-matched no-graph head with the same tuning budget. \\
duplicate items & checked ok & [V] Shipped splits (Laptop 2328/638, Restaurants 3608/1120 aspects, matching Dai et al.'s statistics): one exact duplicate sentence per dataset, same label; an aspect-term vote reaches 52.5\% vs majority 53.4\% (Laptop) and 65.4\% vs 65.0\% (Restaurants), so no aspect-identity shortcut. \cite{dai2021syntax} & not reported \\
train test overlap & checked ok & [V] Shipped splits (Laptop 2328/638, Restaurants 3608/1120 aspects, matching Dai et al.'s statistics): one exact duplicate sentence per dataset, same label; an aspect-term vote reaches 52.5\% vs majority 53.4\% (Laptop) and 65.4\% vs 65.0\% (Restaurants), so no aspect-identity shortcut. \cite{dai2021syntax} & not reported \\
pretraining contamination & not checked & SemEval-2014 reviews have been public since 2014 and may be in RoBERTa/DeBERTa pre-training data; this affects both arms of a with/without-syntax comparison alike. & not reported \\
\bottomrule
\end{longtable}
}

\textbf{Decision: CONDITIONAL GO.} Feasible and useful as a replication; novelty is low by design.

\textbf{Conditions before committing.}
\begin{itemize}
\item Read Lebmeier et al. through the library; if it already re-runs a syntax-GCN model with validation splits and uncertainty, narrow to the DeBERTa-v3 part or drop D1.
\item Fix the validation protocol, seeds and the smallest effect of interest (e.g. 1 point) before any run.
\end{itemize}

\textbf{Unresolved evidence and risks.}
\begin{itemize}
\item Lebmeier et al. unread
\item Reimplementation may not match Dai et al.'s unreleased RoBERTa+tree code
\end{itemize}

\subsection{2. D2 (extension): is it syntax, or just an extra graph head? Structure controls on a fixed fine-tuned encoder}
\textbf{Type of gap:} Underexplored experimental setting\\
\textbf{Recommended mode:} extension

\textbf{Research question.} With the encoder and head parameter count fixed, do dependency trees beat left/right-chain, degree-matched random trees and a fully connected graph inside the same GCN head?

\textbf{Motivation.} Dai et al. ran chain-tree controls only with GloVe models \cite{dai2021syntax}; random masks also help in Syntax-BERT \cite{bai2021syntaxbert}; BERT4GCN's graph adds only slightly \cite{xiao2021bert4gcn}.

\textbf{Supporting papers.} \cite{dai2021syntax,bai2021syntaxbert,xiao2021bert4gcn}

\textbf{Closest existing work.}
\begin{itemize}
\item \cite{dai2021syntax} [read: full text; overlap: partial]: Chain-tree controls with GloVe embeddings only, 3 runs.
\item \cite{bai2021syntaxbert} [read: full text; overlap: partial]: Random-mask control inside the transformer, GLUE tasks, not ABSA.
\item \cite{xiao2021bert4gcn} [read: full text; overlap: partial]: Removes the dependency graph but has no random or chain-tree control.
\item \cite{lebmeier2023repro} [read: abstract; overlap: unknown]: Unknown architectures; unlikely to include structure controls.
\end{itemize}

\textbf{Evidence for the gap.}
\begin{itemize}
\item [V] Chain-tree controls in Dai et al. use 300-d GloVe embeddings kept fixed, not a pretrained encoder. \cite{dai2021syntax}
\item [V] A random decomposition of the network also gives moderate improvements (Syntax-BERT, GLUE). \cite{bai2021syntaxbert}
\end{itemize}

\textbf{Potential contribution.} Separates the effect of syntax from that of adding any graph-structured head.

\textbf{Novelty confidence:} low (scale: very\_low, low, medium; a search cannot establish more). No ABSA paper with random/chain-tree controls on a pretrained encoder was found in the full texts read or by web search; ablation tables of the many 2022-2026 GCN papers were not all read and may contain such controls.
Attempts to disprove novelty: 10 search(es); signal \texttt{no\_close\_match\_in\_this\_search}. Queries: \texttt{does syntax help pretrained transformers aspect-based sentiment}; \texttt{dependency tree graph neural network on top of BERT ablation sentiment}; \texttt{syntactic information unnecessary strong RoBERTa baseline aspect sentiment}; \texttt{syntax-agnostic versus syntax-aware models frozen encoder probing}; \texttt{random trees versus dependency trees graph convolution evaluation}; \texttt{web: does syntax help pretrained transformers dependency tree GNN ablation aspect-based sentiment random trees}; \texttt{web: aspect-based sentiment GCN reproducibility seeds variance same encoder syntax benefit}; \texttt{web: aspect sentiment random tree / fully connected graph replace dependency tree GCN ablation BERT}

\textbf{Feasibility.} [E] Same pipeline as D1 plus graph generators; adds about 40 runs.

\textbf{Methodological risks.}

{\small
\begin{longtable}{>{\raggedright\arraybackslash}p{0.23\textwidth}>{\raggedright\arraybackslash}p{0.23\textwidth}>{\raggedright\arraybackslash}p{0.23\textwidth}>{\raggedright\arraybackslash}p{0.23\textwidth}}
\toprule
\textbf{Risk} & \textbf{Status} & \textbf{Evidence} & \textbf{Mitigation} \\
\midrule
\endhead
model selection on test & found & [V] Official code of the closest studies selects on the test set: RoBERTaABSA finetune.py passes the test split as dev\_data and reloads the best model; ASGCN (original and Dai et al.'s copy) keeps max\_test\_acc and max\_test\_f1 separately, so accuracy and F1 can come from different epochs; RGAT run.py reports the best evaluation on the test set. \cite{dai2021syntax,zhang2019asgcn,wang2020rgat} & Validation split carved from training, grouped by sentence; select checkpoints and hyperparameters on it only; touch the test set once. \\
metric definition & found & [V] ASGCN reports maximum accuracy and maximum F1 from possibly different epochs. \cite{zhang2019asgcn} & Report accuracy and macro-F1 of one selected checkpoint. \\
insufficient runs & found & [V] Dai et al.: 3 runs for the tree comparison, no variance for the RoBERTa table, released logs single-run; Lebmeier et al. report notable variance across data splits. \cite{dai2021syntax,lebmeier2023repro} & At least 10 seeds x repeated train/validation splits; paired comparisons; confidence intervals. \\
weak baseline & found & [V] The same RoBERTa-MLP baseline is reported at 83.78 (Dai et al.) and 81.11 (BERT4GCN) on Laptop14, so cross-paper syntax gains can be baseline artefacts. \cite{dai2021syntax,xiao2021bert4gcn} & Tune the plain-encoder baseline with the same budget and search space as the tree heads. \\
protocol mismatch & found & [V] No public code for the RoBERTa+tree rows; the Hugging Face SetFit copy has 654 laptop test rows vs 638 aspects in the shipped split. \cite{dai2021syntax} & Use the shipped split files; reimplement heads on one codebase; label the work a replication. \\
confound & not checked & Tree heads add parameters and a separate learning rate; Dai et al. note RoBERTa-based tree models are hard to optimise and need different learning rates. & Parameter-matched no-graph head with the same tuning budget. \\
duplicate items & checked ok & [V] Shipped splits (Laptop 2328/638, Restaurants 3608/1120 aspects, matching Dai et al.'s statistics): one exact duplicate sentence per dataset, same label; an aspect-term vote reaches 52.5\% vs majority 53.4\% (Laptop) and 65.4\% vs 65.0\% (Restaurants), so no aspect-identity shortcut. \cite{dai2021syntax} & not reported \\
train test overlap & checked ok & [V] Shipped splits (Laptop 2328/638, Restaurants 3608/1120 aspects, matching Dai et al.'s statistics): one exact duplicate sentence per dataset, same label; an aspect-term vote reaches 52.5\% vs majority 53.4\% (Laptop) and 65.4\% vs 65.0\% (Restaurants), so no aspect-identity shortcut. \cite{dai2021syntax} & not reported \\
pretraining contamination & not checked & SemEval-2014 reviews have been public since 2014 and may be in RoBERTa/DeBERTa pre-training data; this affects both arms of a with/without-syntax comparison alike. & not reported \\
\bottomrule
\end{longtable}
}

\textbf{Decision: CONDITIONAL GO.} The most informative design: it tests the mechanism, and runs on D1's pipeline.

\textbf{Conditions before committing.}
\begin{itemize}
\item Run D1 first so the baseline and protocol are trusted.
\item Screen ablation tables of the 5-10 most cited 2022-2026 syntax-GCN ABSA papers for random or chain-tree controls on PLMs.
\end{itemize}

\textbf{Unresolved evidence and risks.}
\begin{itemize}
\item Uncited ablations in recent GCN papers
\item Effect sizes may be below the detectable difference on 638/1,120 test aspects
\end{itemize}

\subsection{3. D3 (novel investigation): frozen versus fine-tuned encoder - does syntax matter only when the encoder cannot adapt?}
\textbf{Type of gap:} Unverified novelty hypothesis\\
\textbf{Recommended mode:} original

\textbf{Research question.} When the encoder is frozen (features only), do dependency-tree heads give larger gains over a parameter-matched no-graph head than when it is fine-tuned?

\textbf{Motivation.} Every closest study fine-tunes the encoder \cite{dai2021syntax,sachan2021syntax}; Dai et al. argue fine-tuning makes the encoder learn task-oriented structure implicitly.

\textbf{Supporting papers.} \cite{dai2021syntax,sachan2021syntax}

\textbf{Closest existing work.}
\begin{itemize}
\item \cite{sachan2021syntax} [read: full text; overlap: partial]: Fine-tuned encoders only, information extraction tasks.
\item \cite{dai2021syntax} [read: full text; overlap: partial]: Fine-tuned RoBERTa only; frozen GloVe for the tree comparison.
\item \cite{hou2021sagcn} [read: abstract; overlap: unknown]: BERT + GCN depth variants; encoder regime not read.
\end{itemize}

\textbf{Evidence for the gap.}
\begin{itemize}
\item [V] The closest full texts fine-tune the encoder in every syntax comparison. \cite{dai2021syntax,sachan2021syntax}
\item [S] If fine-tuning lets the encoder absorb task structure, a frozen encoder should leave more room for explicit syntax.
\end{itemize}

\textbf{Potential contribution.} Identifies when explicit syntax is worth adding (compute-limited, frozen-feature settings).

\textbf{Novelty confidence:} low (scale: very\_low, low, medium; a search cannot establish more). Unverified hypothesis: probing and adapter literature with frozen encoders was not searched in depth.
Attempts to disprove novelty: 10 search(es); signal \texttt{no\_close\_match\_in\_this\_search}. Queries: \texttt{does syntax help pretrained transformers aspect-based sentiment}; \texttt{dependency tree graph neural network on top of BERT ablation sentiment}; \texttt{syntactic information unnecessary strong RoBERTa baseline aspect sentiment}; \texttt{syntax-agnostic versus syntax-aware models frozen encoder probing}; \texttt{random trees versus dependency trees graph convolution evaluation}; \texttt{web: does syntax help pretrained transformers dependency tree GNN ablation aspect-based sentiment random trees}; \texttt{web: aspect-based sentiment GCN reproducibility seeds variance same encoder syntax benefit}; \texttt{web: aspect sentiment random tree / fully connected graph replace dependency tree GCN ablation BERT}

\textbf{Feasibility.} [E] Frozen features can be cached once; cheapest of all designs (CPU-trainable heads).

\textbf{Methodological risks.}

{\small
\begin{longtable}{>{\raggedright\arraybackslash}p{0.23\textwidth}>{\raggedright\arraybackslash}p{0.23\textwidth}>{\raggedright\arraybackslash}p{0.23\textwidth}>{\raggedright\arraybackslash}p{0.23\textwidth}}
\toprule
\textbf{Risk} & \textbf{Status} & \textbf{Evidence} & \textbf{Mitigation} \\
\midrule
\endhead
model selection on test & found & [V] Official code of the closest studies selects on the test set: RoBERTaABSA finetune.py passes the test split as dev\_data and reloads the best model; ASGCN (original and Dai et al.'s copy) keeps max\_test\_acc and max\_test\_f1 separately, so accuracy and F1 can come from different epochs; RGAT run.py reports the best evaluation on the test set. \cite{dai2021syntax,zhang2019asgcn,wang2020rgat} & Validation split carved from training, grouped by sentence; select checkpoints and hyperparameters on it only; touch the test set once. \\
metric definition & found & [V] ASGCN reports maximum accuracy and maximum F1 from possibly different epochs. \cite{zhang2019asgcn} & Report accuracy and macro-F1 of one selected checkpoint. \\
insufficient runs & found & [V] Dai et al.: 3 runs for the tree comparison, no variance for the RoBERTa table, released logs single-run; Lebmeier et al. report notable variance across data splits. \cite{dai2021syntax,lebmeier2023repro} & At least 10 seeds x repeated train/validation splits; paired comparisons; confidence intervals. \\
weak baseline & found & [V] The same RoBERTa-MLP baseline is reported at 83.78 (Dai et al.) and 81.11 (BERT4GCN) on Laptop14, so cross-paper syntax gains can be baseline artefacts. \cite{dai2021syntax,xiao2021bert4gcn} & Tune the plain-encoder baseline with the same budget and search space as the tree heads. \\
protocol mismatch & found & [V] No public code for the RoBERTa+tree rows; the Hugging Face SetFit copy has 654 laptop test rows vs 638 aspects in the shipped split. \cite{dai2021syntax} & Use the shipped split files; reimplement heads on one codebase; label the work a replication. \\
confound & not checked & Tree heads add parameters and a separate learning rate; Dai et al. note RoBERTa-based tree models are hard to optimise and need different learning rates. & Parameter-matched no-graph head with the same tuning budget. \\
duplicate items & checked ok & [V] Shipped splits (Laptop 2328/638, Restaurants 3608/1120 aspects, matching Dai et al.'s statistics): one exact duplicate sentence per dataset, same label; an aspect-term vote reaches 52.5\% vs majority 53.4\% (Laptop) and 65.4\% vs 65.0\% (Restaurants), so no aspect-identity shortcut. \cite{dai2021syntax} & not reported \\
train test overlap & checked ok & [V] Shipped splits (Laptop 2328/638, Restaurants 3608/1120 aspects, matching Dai et al.'s statistics): one exact duplicate sentence per dataset, same label; an aspect-term vote reaches 52.5\% vs majority 53.4\% (Laptop) and 65.4\% vs 65.0\% (Restaurants), so no aspect-identity shortcut. \cite{dai2021syntax} & not reported \\
pretraining contamination & not checked & SemEval-2014 reviews have been public since 2014 and may be in RoBERTa/DeBERTa pre-training data; this affects both arms of a with/without-syntax comparison alike. & not reported \\
\bottomrule
\end{longtable}
}

\textbf{Decision: CONDITIONAL GO.} Plausible and cheap, but the novelty rests on absence of evidence.

\textbf{Conditions before committing.} Run a dedicated challenge and web search on frozen-encoder syntax probing (e.g. 'frozen BERT features dependency GCN') before committing.

\textbf{Unresolved evidence and risks.} Probing literature may already answer it

\subsection{Candidates rejected after the novelty challenge}

\begin{itemize}
\item \textbf{As stated: do syntax-tree GCNs help when the pretrained encoder is held fixed?} (Already well studied): Already performed: Dai et al. compare RoBERTa-MLP with RoBERTa+ASGCN/PWCN/RGAT using parser and induced trees and conclude that trees bring no significant improvement \cite{dai2021syntax}; BERT4GCN's own ablation finds the dependency graph adds only slightly \cite{xiao2021bert4gcn}; and the same comparison under identical encoders is done in information extraction \cite{sachan2021syntax}.
\item \textbf{D4 (extension): parser quality as the moderator of syntax gains in ABSA} (Underexplored experimental setting): Sachan et al. show gains depend on gold parses \cite{sachan2021syntax}; SemEval reviews have no gold dependency trees, so parse quality cannot be measured against a reference and the moderator cannot be isolated.
\end{itemize}

\section{Feasibility ranking}

Each criterion is scored 0-5 with evidence. score\_percent is the weighted share earned on scored criteria only; score\_range shows the total if unscored criteria were 0 or 5. Any blocker removes a candidate from the ranking no matter how high its other scores are. Novelty is capped at 2 unless the closest prior work was read in full, and at 1 when that work already ran the experiment.

{\small
\begin{longtable}{>{\raggedright\arraybackslash}p{0.131\textwidth}>{\raggedright\arraybackslash}p{0.131\textwidth}>{\raggedright\arraybackslash}p{0.131\textwidth}>{\raggedright\arraybackslash}p{0.131\textwidth}>{\raggedright\arraybackslash}p{0.131\textwidth}>{\raggedright\arraybackslash}p{0.131\textwidth}>{\raggedright\arraybackslash}p{0.131\textwidth}}
\toprule
\textbf{Rank} & \textbf{Candidate} & \textbf{Status} & \textbf{Score \%} & \textbf{Range} & \textbf{Scored weight \%} & \textbf{Verified weight \%} \\
\midrule
\endhead
1 & D1 replication: protocol-corrected re-test & VIABLE & 71.8 & 71.8-71.8 & 100.0 & 61.0 \\
2 & D3 novel: frozen vs fine-tuned encoder & VIABLE & 68.2 & 68.2-68.2 & 100.0 & 30.0 \\
3 & D2 extension: structure controls & VIABLE & 64.6 & 64.6-64.6 & 100.0 & 37.0 \\
4 & D4 parser-quality moderator & VIABLE & 53.8 & 53.8-53.8 & 100.0 & 49.0 \\
\bottomrule
\end{longtable}
}

\textbf{Scores per criterion (0-5).}

{\small
\begin{longtable}{>{\raggedright\arraybackslash}p{0.153\textwidth}>{\raggedright\arraybackslash}p{0.153\textwidth}>{\raggedright\arraybackslash}p{0.153\textwidth}>{\raggedright\arraybackslash}p{0.153\textwidth}>{\raggedright\arraybackslash}p{0.153\textwidth}>{\raggedright\arraybackslash}p{0.153\textwidth}}
\toprule
\textbf{Criterion} & \textbf{Weight} & \textbf{D1} & \textbf{D3} & \textbf{D2} & \textbf{D4} \\
\midrule
\endhead
scientific relevance & 12 & 4 [V] & 3 [S] & 4 [E] & 2 [V] \\
novelty or replication value & 12 & 4 [V] & 3 [E] & 3 [E] & 2 [E] \\
paper and code availability & 10 & 3 [V] & 3 [V] & 3 [V] & 3 [V] \\
dataset accessibility & 12 & 3 [V] & 3 [V] & 3 [V] & 3 [V] \\
compute fit & 12 & 4 [E] & 4 [E] & 4 [E] & 4 [E] \\
implementation simplicity & 10 & 3 [E] & 4 [E] & 3 [E] & 3 [E] \\
time fit & 10 & 4 [E] & 5 [E] & 3 [E] & 4 [E] \\
reproducibility & 8 & 2 [V] & 2 [V] & 2 [V] & 2 [V] \\
evaluation clarity & 7 & 4 [V] & 3 [E] & 3 [V] & 1 [V] \\
low failure risk & 7 & 5 [E] & 4 [E] & 4 [E] & 2 [E] \\
\bottomrule
\end{longtable}
}

\textbf{D1 replication: protocol-corrected re-test}
\begin{itemize}
\item scientific relevance: Null result is cited as settled; claims of gains continue \cite{li2024sebert}
\item novelty or replication value: Single-table, test-selected result that many papers rely on \cite{dai2021syntax}
\item paper and code availability: Papers open; tree-head code GloVe-only, pinned to 2020 libraries
\item dataset accessibility: Shipped splits with parses; counts match the paper
\item compute fit: Base encoders, \textasciitilde{}3.6k training aspects
\item implementation simplicity: Reimplement heads on current transformers
\item time fit: About 4 weeks
\item reproducibility: Test-set selection, few runs, missing RoBERTa+tree code
\item evaluation clarity: Standard metrics, published numbers to compare
\item low failure risk: Either outcome reportable
\item Missing: dataset\_accessibility: No licence
\end{itemize}

\textbf{D3 novel: frozen vs fine-tuned encoder}
\begin{itemize}
\item scientific relevance: Relevant to compute-limited practice
\item novelty or replication value: Not found in the closest full texts
\item paper and code availability: Papers open; tree-head code GloVe-only, pinned to 2020 libraries
\item dataset accessibility: Shipped splits with parses; counts match the paper
\item compute fit: Base encoders, \textasciitilde{}3.6k training aspects
\item implementation simplicity: Cached features, small heads
\item time fit: About 4 weeks
\item reproducibility: Test-set selection, few runs, missing RoBERTa+tree code
\item evaluation clarity: No published numbers
\item low failure risk: Either outcome reportable
\item Missing: novelty\_or\_replication\_value: Probing literature unsearched
\item Missing: dataset\_accessibility: No licence
\end{itemize}

\textbf{D2 extension: structure controls}
\begin{itemize}
\item scientific relevance: Tests the mechanism behind a large literature
\item novelty or replication value: No PLM structure controls found in the full texts read
\item paper and code availability: Papers open; tree-head code GloVe-only, pinned to 2020 libraries
\item dataset accessibility: Shipped splits with parses; counts match the paper
\item compute fit: Base encoders, \textasciitilde{}3.6k training aspects
\item implementation simplicity: Graph generators on D1 pipeline
\item time fit: 6 weeks including D1
\item reproducibility: Test-set selection, few runs, missing RoBERTa+tree code
\item evaluation clarity: No published numbers for the controls
\item low failure risk: Null or positive both informative
\item Missing: novelty\_or\_replication\_value: Recent ablation tables unread
\item Missing: dataset\_accessibility: No licence
\end{itemize}

\textbf{D4 parser-quality moderator}
\begin{itemize}
\item scientific relevance: Shown for IE \cite{sachan2021syntax}
\item novelty or replication value: Moderator not isolable without gold parses
\item paper and code availability: Papers open; tree-head code GloVe-only, pinned to 2020 libraries
\item dataset accessibility: Shipped splits with parses; counts match the paper
\item compute fit: Base encoders, \textasciitilde{}3.6k training aspects
\item implementation simplicity: Multiple parsers
\item time fit: 5 weeks
\item reproducibility: Test-set selection, few runs, missing RoBERTa+tree code
\item evaluation clarity: No gold trees for review data
\item low failure risk: Likely inconclusive
\item Missing: dataset\_accessibility: No licence
\end{itemize}

\section{Candidate datasets}

{\small
\begin{longtable}{>{\raggedright\arraybackslash}p{0.102\textwidth}>{\raggedright\arraybackslash}p{0.102\textwidth}>{\raggedright\arraybackslash}p{0.102\textwidth}>{\raggedright\arraybackslash}p{0.102\textwidth}>{\raggedright\arraybackslash}p{0.102\textwidth}>{\raggedright\arraybackslash}p{0.102\textwidth}>{\raggedright\arraybackslash}p{0.102\textwidth}>{\raggedright\arraybackslash}p{0.102\textwidth}>{\raggedright\arraybackslash}p{0.102\textwidth}}
\toprule
\textbf{Dataset} & \textbf{Purpose} & \textbf{Size} & \textbf{Format} & \textbf{Licence} & \textbf{Access} & \textbf{Preprocessing} & \textbf{Suitability} & \textbf{Verified via API} \\
\midrule
\endhead
\href{https://github.com/ROGERDJQ/RoBERTaABSA/tree/main/Dataset}{SemEval-2014 Laptop and Restaurant splits shipped with RoBERTaABSA} & All designs (reference split) & Laptop 2,328/638 aspects; Restaurants 3,608/1,120 (counted from Train/Test.json) & JSON with tokens, POS, dependency heads and relations; original XML & not reported & reachable (HTTP 200) & Conflict polarities removed; parses precomputed & Matches Dai et al.'s statistics exactly; repository has no licence & yes \\
\href{https://huggingface.co/datasets/tomaarsen/setfit-absa-semeval-restaurants}{SemEval-2014 restaurants (SetFit copy)} & Not recommended for comparison & 3,693 / 1,134 rows & parquet & not reported & open, not gated & not reported & Row counts differ from the published split (laptop copy: 654 test rows vs 638 aspects) & yes \\
\href{https://huggingface.co/datasets/tomaarsen/setfit-absa-semeval-laptops}{SemEval-2014 laptops (SetFit copy)} & Not recommended for comparison & 2,358 / 654 rows & parquet & not reported & open, not gated & not reported & As above & yes \\
\bottomrule
\end{longtable}
}

\section{Recommendation}

\begin{itemize}
\item [E] No-go for the question as stated, and for the parser-quality variant.
\item [E] Conditional go for D1 followed by D2 as one 6-8 week project: replicate the RoBERTa tree-head comparison with validation-based selection and 10+ seeds, then add chain, random and fully connected graph controls with matched parameters.
\item [E] First action, no compute: read Lebmeier et al. through the library and screen ablation tables of the most cited 2022-2026 syntax-GCN papers for structure controls.
\item [S] D3 (frozen encoder) is the only design with a novel framing; check the probing literature first.
\end{itemize}

\section{Limitations of this report}

\begin{itemize}
\item Lebmeier et al. (2023) was not read beyond its abstract; it may already contain part of D1.
\item The ablation tables of most 2022-2026 syntax-GCN ABSA papers were not read; D2's controls may exist there.
\item Semantic Scholar was rate-limited; the challenge command did not retrieve three of the closest papers, which were found by web search.
\item Compute and time figures are estimates; no model was trained and no code was run beyond data checks.
\item SemEval-2014 data licence is not stated on any copy checked.
\end{itemize}

\section{References}

\bibliographystyle{plain}
\bibliography{references}

\end{document}
